\documentclass[11pt]{article}
\usepackage[margin=1in]{geometry}
\usepackage{amsmath,amssymb,amsthm,booktabs}
\usepackage[hidelinks]{hyperref}
\newtheorem{theorem}{Theorem}
\newtheorem{proposition}{Proposition}
\newtheorem{corollary}{Corollary}
\newcommand{\norm}[1]{\left\lVert#1\right\rVert}
\title{Adaptation Without Strict Execution Contraction\\
\large A proposed theory section for frozen VLA policies}
\author{}
\date{September 14, 2026}
\begin{document}
\maketitle

\noindent
This draft synthesizes the evidence at repository commit \texttt{b540d52}.
Its mathematical statements are conditional; the experiments do not certify
their physical bounds. The supporting note records the empirical evidence,
counterexamples, and a corrected telemetry audit. This standalone draft is
not inserted into the submission manuscript.

\section{Observation, correction, and task execution}
A frozen policy produces a nominal command $a_k$, and the adapter sends
$u_k=a_k-c_k$. The healthy predictor must use the known command $u_k$.
At a specified state and horizon, distinguish the residual observation
$r_k$ from the physical discrepancy $v_k$ remaining after correction:
\begin{equation}
 r_k=H_k\theta_k+\varepsilon_k,\qquad
 v_k=F_k\theta_k-D_kc_k.
 \label{eq:maps}
\end{equation}
The physical disturbance map $F_k$, observation map $H_k$, and available-input
map $D_k$ need not coincide. For any nominal allocation $K_k$,
\begin{equation}
 v_k=(F_k-D_kK_k)\theta_k+
 D_kK_k(\theta_k-\hat\theta_k)+D_k(K_k\hat\theta_k-c_k).
 \label{eq:decomposition}
\end{equation}
The identity separates error in the cancellation map, estimation error,
and error in allocating and realizing the correction.
The first term concerns the chosen allocation; it need not be irreducible.

In the noiseless unconstrained linear case, a linear rule $c=Lr$ cancels
every disturbance exactly iff
\begin{equation}
 \ker H\subseteq\ker F,\qquad
 \operatorname{range}F\subseteq\operatorname{range}D.
\end{equation}
Indeed, the first inclusion makes $T(H\theta)=F\theta$ well-defined and the
second permits lifting $T$ through $D$. Constraints additionally require
pointwise feasibility $F\theta\in D\mathcal C$. Existence of one feasible
linear rule over a disturbance set further requires $LH\theta\in\mathcal C$
throughout that set. Accurate observation therefore does not imply feasible
cancellation, and a reduced fitted objective does not by itself imply
physical improvement.

\section{Two implemented estimator recursions}
For command-equivalent additive faults, write
\[
 z_k=M^{-1}r_k-b=f_k+w_k,\quad e_k=\hat f_k-f_k,\quad E_k=\norm{e_k},
 \quad \nu_k=\norm{f_{k+1}-f_k}.
\]
Assume $f_k\in\mathcal C$, with $\mathcal C$ closed and convex, Euclidean
projection $\Pi_{\mathcal C}$, and $\norm{w_k}\le\epsilon_k$.
The latter must be independently bounded to obtain a predictive guarantee.

\begin{proposition}[Legacy and innovation bounds]
For $0<\gamma\le1$ and $0\le s_k\le1$, the legacy update
\[
\hat f_{k+1}=\Pi_{\mathcal C}
\big[(1-\gamma)\hat f_k+\gamma s_kz_k\big]
\]
satisfies
\begin{equation}
E_{k+1}\le(1-\gamma)E_k+\gamma s_k\epsilon_k+
              \gamma(1-s_k)\norm{f_k}+\nu_k.
\label{eq:legacy}
\end{equation}
For $0\le\alpha_k\le1$, the innovation update
\[
\hat f_{k+1}=\Pi_{\mathcal C}
\big[\hat f_k+\alpha_k(z_k-\hat f_k)\big]
\]
satisfies
\begin{equation}
E_{k+1}\le(1-\alpha_k)E_k+\alpha_k\epsilon_k+\nu_k.
\label{eq:innovation}
\end{equation}
\end{proposition}
\begin{proof}
Before projection, subtract $f_k$. The respective expressions are
$(1-\gamma)e_k+\gamma s_kw_k-\gamma(1-s_k)f_k$ and
$(1-\alpha_k)e_k+\alpha_kw_k$.
Nonexpansiveness of projection relative to $f_k$, the triangle inequality,
and the change from $f_k$ to $f_{k+1}$ give the results.
\end{proof}
The legacy zero-target deadzone is $s_k=0$; an explicit hold branch is a
different update. Innovation uses $\alpha_k=\gamma s_k$, with its normalizer
evaluated on the innovation. Persistent $w_k$ biases both estimators.
Only legacy has the additional attenuation forcing in
\eqref{eq:legacy}. For constant positive $\alpha$, the innovation steady
bound is $\epsilon+\nu/\alpha$.
Holding at $T$ instead gives
\begin{equation}
\norm{\hat f_T-f_{T+j}}\le E_T+
      \sum_{\ell=T}^{T+j-1}\nu_\ell.
\label{eq:hold}
\end{equation}
This also describes snapshot age. Holding removes subsequent estimator
updates; it retains calibration error and fault-drift sensitivity.

\section{Finite-horizon incremental storage}
Augment the physical deviation with the controller/decoder memory,
predictor history, estimator state, and relevant policy/chunk state.
Let the resulting deviation be $Z_k$. Relative to a specified nominal
trajectory, suppose the incremental dynamics on a declared domain satisfy
\begin{equation}
 Z_{k+1}=\mathcal A_kZ_k+\mathcal B_kd_k+\rho_k.
 \label{eq:augmented}
\end{equation}
The signed transition includes estimator--execution feedback; $d_k,\rho_k$
include disturbance drift, unmatched effects, nonlinear remainder, and
unmodeled inputs. A healthy output-prediction fit alone does not establish
\eqref{eq:augmented} or uniform remainder bounds.

\begin{theorem}[Finite-horizon storage bound]
Fix a finite horizon $N$. Choose a positive-definite metric $P_k$ at each
time $k=0,\ldots,N$. Define
\[
 V_k=Z_k^\top P_kZ_k,\qquad
 a_k=\norm{P_{k+1}^{1/2}\mathcal A_kP_k^{-1/2}}_2,
\]
and suppose
$\norm{\mathcal B_kd_k+\rho_k}_{P_{k+1}}\le b_k$.
If $R_0\ge\sqrt{V_0}$, then
\begin{align}
 \sqrt{V_k}&\le R_k,\qquad R_{k+1}=a_kR_k+b_k,\\
 R_n&=\left(\prod_{j=0}^{n-1}a_j\right)R_0+
 \sum_{k=0}^{n-1}\left(\prod_{j=k+1}^{n-1}a_j\right)b_k.
 \label{eq:tube}
\end{align}
\end{theorem}
\begin{proof}
The induced-norm and triangle inequalities give
$\sqrt{V_{k+1}}\le a_k\sqrt{V_k}+b_k$.
Induction yields both claims.
\end{proof}
The theorem permits $a_k=1$ and finite periods with $a_k>1$.
It establishes a finite-horizon bound, not asymptotic stability.
Report metric coercivity and physical-output conversion: shrinking $P_k$
does not make a growing physical deviation harmless.
For local assumptions, verify the tube stays within their domain.
Keeping signed matrix transition products can be substantially tighter than
products of scalar norms.

\begin{corollary}[Strict contraction as a special case]
If $a_k\le a<1$ and $b_k\le b$, then
$R_n\le a^nR_0+b(1-a^n)/(1-a)$.
\end{corollary}
Alternatively, a separately justified physical inequality
$X_{k+1}\le\lambda_kX_k+L_k\norm{v_k}+\eta_k$
fits the theorem with \eqref{eq:decomposition}. Its input is the remaining
physical discrepancy $v_k$, not the predictor residual $r_k$.
Contact/controller resets require additional jump inequalities and domain
conditions; a fitted free-space metric does not cover them automatically.

\section{Execution memory explains why identification and repair differ}
Consider an exact scalar example, $0<\alpha\le1$:
\[
p_{k+1}=\lambda p_k+g(f-\hat f_k),\qquad
\hat f_{k+1}=(1-\alpha)\hat f_k+\alpha f,\qquad p_0=\hat f_0=0.
\]
The remaining fault is $f(1-\alpha)^k$. For $\lambda\ne1-\alpha$,
\begin{equation}
p_n=gf\frac{\lambda^n-(1-\alpha)^n}{\lambda-(1-\alpha)};
\end{equation}
the coincident-root value is $gf\,n\lambda^{n-1}$.
For a stable servo, $0\le\lambda<1$, the displacement decays. For an
integrator, $\lambda=1$,
\begin{equation}
p_n=\frac{gf}{\alpha}[1-(1-\alpha)^n]\longrightarrow\frac{gf}{\alpha}.
\label{eq:transient-cost}
\end{equation}
Thus an unbiased observer stops further drift but need not erase transient
displacement. With legacy attenuation $0\le s<1$ and $0<\gamma\le1$,
\begin{equation}
p_n=gf\left[(1-s)n+\frac{s}{\gamma}
                   (1-(1-\gamma)^n)\right].
\end{equation}
Persistent bias and transient cost enter differently. For an ideal
integrator held for $h$ steps, \eqref{eq:hold} bounds the subsequent
displacement contribution by
$|g|[hE_T+\nu h(h-1)/2]$.
Identification followed by a task reset can remove the earlier displacement;
merely pausing an estimate during a disrupted trajectory does not.
These are constructed examples, not fitted explanations of individual runs.

\section{Joint Lyapunov stability without physical subsystem contraction}
Consider a normalized neutral channel, constant fault, exact observation
$z_k=f$, and no projection. Apply the composite update
$\hat f_{k+1}=\hat f_k+\alpha(z_k-\hat f_k)+\kappa e_{k+1}$.
With $\delta_k=\hat f_k-f$, the error system is
\begin{equation}
 e_{k+1}=e_k-\delta_k,\qquad
 \delta_{k+1}=(1-\alpha)\delta_k+\kappa e_{k+1}.
\end{equation}
\begin{proposition}[Stabilization of a neutral execution mode]
The joint matrix
\[
\mathcal A=\begin{bmatrix}1&-1\\
                   \kappa&1-\alpha-\kappa\end{bmatrix}
\]
is Schur stable iff $0<\alpha<2$ and $0<\kappa<4-2\alpha$.
\end{proposition}
\begin{proof}
$\det\mathcal A=1-\alpha$ and
$\operatorname{tr}\mathcal A=2-\alpha-\kappa$.
The strict Jury inequalities reduce to
$\alpha>0$, $\kappa>0$, and $4-2\alpha-\kappa>0$.
\end{proof}
For every $Q\succ0$ there consequently exists $P\succ0$ satisfying
$\mathcal A^\top P\mathcal A-P=-Q$.
For $Z_{k+1}=\mathcal AZ_k+w_k$ and $q=\lambda_{\min}(Q)$,
\begin{equation}
V_{k+1}-V_k\le-\frac q2\norm{Z_k}^2+
\left(\norm{P}_2+
\frac{2\norm{\mathcal A^\top P}_2^2}{q}\right)\norm{w_k}^2.
\end{equation}
This follows by expanding the quadratic and applying Young's inequality.
It gives a joint Lyapunov/ISS bound even though the physical subsystem
alone is neutral. At $\kappa=0$ the neutral eigenvalue remains; excessive
tracking gain destabilizes the system.

This is a conditional design result. The actual reference, covariance,
input map, sampling, clipping, and contact behavior must be included before
applying it to a robot. Increment-only tracking may not observe a neutral
pose mode. Stable individual frozen-gain matrices do not establish stability
under arbitrary gain switching.

\section{Conditional task transfer and empirical interpretation}
Same-command healthy physical execution differs from the healthy closed-loop
policy counterfactual. A frozen policy generally emits different commands
after a disturbance. Include that feedback in $Z_k$, or bound the command
difference separately; a predictor residual is neither trajectory.

Suppose a successful healthy trajectory has robustness radius $\mu>0$ in
the weighted supremum metric
$d_{\rm task}=\max_{k\le N}\norm{W_k(x_k-x_k^0)}$.
If $x_k-x_k^0=C_kZ_k$ and
\begin{equation}
\max_{k\le N}\norm{W_kC_kP_k^{-1/2}}_2R_k<\mu,
\end{equation}
with compatible domain and hybrid assumptions, the corrected trajectory
also succeeds. This is a consequence of the assumed success neighborhood,
not a margin identified by the present experiments. Static-offset sweeps
do not characterize arbitrary time-varying trajectory errors.

The evidence motivates this distinction: same-command Panda perturbations
do not establish strict contraction, while repair remains observed; a
matched-estimate ALOHA innovation observer repairs without holding;
improved ARX prediction does not remove rotation-y observation bias; and
the pinned static-observer comparison does not resolve a FIR-memory benefit.
The finite-horizon theorem accommodates these outcomes without asserting
they certify its constants. The authority and joint Lyapunov analyses also
permit harmful correction and destabilizing tracking feedback.

The theoretical tools are standard. The proposed contribution is their
faithful connection to the implemented VLA interface and explicit empirical
tests of the required assumptions, not universal adaptive recovery.

\begin{thebibliography}{9}
\bibitem{forni2014}
F. Forni and R. Sepulchre.
A differential Lyapunov framework for contraction analysis.
\emph{IEEE Transactions on Automatic Control}, 2014.
\url{https://arxiv.org/abs/1208.2943}.
\bibitem{forni2013}
F. Forni and R. Sepulchre.
On differentially dissipative dynamical systems, 2013.
\url{https://arxiv.org/abs/1305.3456}.
\bibitem{geiselhart}
R. Geiselhart and F. Wirth.
Relaxed ISS Small-Gain Theorems for Discrete-Time Systems.
\url{https://arxiv.org/abs/1406.3224}.
\bibitem{cheng}
Y. Cheng et al.
Improving the Robustness of Reinforcement Learning Policies with
$\mathcal L_1$ Adaptive Control.
\emph{IEEE Robotics and Automation Letters}, 2022.
\url{https://arxiv.org/abs/2112.01953}.
\bibitem{slotine}
J.-J. E. Slotine and W. Li.
Composite adaptive control of robot manipulators.
\emph{Automatica}, 1989.
\url{https://doi.org/10.1016/0005-1098(89)90094-0}.
\end{thebibliography}
The nonstrict/incremental and input-supply interpretations follow
\cite{forni2014,forni2013}; finite-step interconnection analysis is discussed
in \cite{geiselhart}. Frozen-policy augmentation and composite adaptation
have established antecedents \cite{cheng,slotine}. These works do not
automatically certify the particular discrete controller studied here.
\end{document}
